\appendix
\section{平均と期待値}

    \cref{thm:mean-expectation}の「ランダムに取り出す」とは, $n$ 個ある位置のどれもが同じ確率 $1/n$ で選ばれる, という意味である.

    まず位置で数えてみる.
    $i$ 番目の位置が選ばれる確率は $1/n$ で, そのとき $X$ は $x_i$ になる.
    期待値は値にその確率をかけて足したものなので,
    \[
      E[X] = \sum_{i=1}^{n} x_i \cdot \frac{1}{n} = \frac{1}{n}\sum_{i=1}^{n} x_i = \bar{x}
    \]
    となる.

    少し気になるのは, 同じ値がいくつも並んでいる場合だ.
    期待値の定義では, 異なる値ごとに「値 × その値が出る確率」を足すのだった.
    そこで, 異なる値 $v$ のそれぞれについて, $v$ が並んでいる位置の個数を $c_v$ とする.
    $X = v$ になるのは $c_v$ 個の位置のどれかが選ばれたときなので, $P(X = v) = c_v/n$ である.
    よって
    \[
      E[X] = \sum_{v} v \cdot \frac{c_v}{n} = \frac{1}{n}\sum_{v} c_v v
    \]
    となる.
    $c_v v$ は同じ値 $v$ を $c_v$ 回足したものだから, 右辺は $x_1$ から $x_n$ までを同じ値どうしでまとめてから足しているだけである.
    まとめ方を変えても合計は変わらないので, これも $\frac{1}{n}\sum_{i=1}^{n} x_i$ に等しく, 位置で数えたときと一致する.

    商品Bの20件で確かめてみる.
    1点が5件, 5点が15件なので, $P(X=1) = 5/20$, $P(X=5) = 15/20$ である.
    \[
      E[X] = 1 \cdot \frac{5}{20} + 5 \cdot \frac{15}{20} = \frac{5 + 75}{20} = 4.0
    \]
    となり, 本文で求めた平均4.0と一致する.
